% ECONOMICS_PROBLEM: {"id": "J16-592", "title": "Child penalties and the gender gap", "jel_code": "J16", "source_id": "OP-0029"}
% !TeX program = xelatex
\documentclass[11pt,a4paper]{article}
\usepackage[margin=23mm]{geometry}
\usepackage{fontspec}
\setmainfont{Latin Modern Roman}
\usepackage{amsmath,amssymb,mathtools}
\usepackage{xcolor,enumitem,booktabs,longtable,array,xurl,hyperref}
\definecolor{muted}{HTML}{53636C}
\hypersetup{colorlinks=true,urlcolor=blue,linkcolor=blue}
\setlength{\parindent}{0pt}
\setlength{\parskip}{5pt}
\newcommand{\R}{\mathbb R}
\newcommand{\E}{\mathbb E}
\newcommand{\N}{\mathbb N}
\newcommand{\ind}{\mathbf 1}
\newcommand{\Var}{\operatorname{Var}}
\newcommand{\Cov}{\operatorname{Cov}}
\newcommand{\supp}{\operatorname{supp}}
\DeclareMathOperator*{\argmin}{arg\,min}
\DeclareMathOperator*{\argmax}{arg\,max}
\newcommand{\entrymeta}[3]{{\sffamily\small\color{muted}\textbf{\texttt{#1}}\quad|\quad #2\par #3\par}}
\newcommand{\Problemref}[1]{\texttt{#1}}
\newcommand{\JELref}[2]{\texttt{#2} (JEL \texttt{#1})}
\begin{document}
\section*{Child penalties and the gender gap}
\textbf{Problem ID:} \texttt{J16-592}\quad \textbf{JEL:} \texttt{J16}\par
% BEGIN ECONOMICS STATEMENT
\label{problem:OP-0029}
\label{atlas:prob:L9}
\noindent\textbf{Original atlas ID:} L9 (entry 29). \textbf{Classification:} Editorial birth-timing and family-policy identification problem.

\noindent\textbf{Original atlas question:} Which combination of household specialization, norms, employer incentives and childcare institutions generates persistent post-childbirth earnings gaps?

\subsection*{Definitions and assumptions}
Fix a baseline cohort, an intended first-birth date $T$, and horizon $H$. Let $g\in\{w,m\}$ label the two recorded parental groups under study; this indexing does not assert that all families have this composition. Define $Y_{i,T+\tau}(b,\pi)$ as earnings under birth at $T$ ($b=1$) or a specified delay beyond $T+H$ ($b=0$), and family/workplace policy $\pi$. Use earnings levels, allowing zero, rather than undefined log earnings for nonworkers. The birth intervention must specify timing and associated fertility choices. Assume well-defined potential outcomes and integrability; causal identification requires an explicit assignment or counterfactual-trend design.
\subsection*{Formal research question}
For event time $\tau\in\{0,\ldots,H\}$, define
\[
 \delta_g(\tau,\pi)=\mathbb E[Y_{i,T+\tau}(1,\pi)-Y_{i,T+\tau}(0,\pi)\mid g],
 \qquad
 \Pi_\tau(\pi)=\frac{\delta_m(\tau,\pi)}{\mu_m}
                 -\frac{\delta_w(\tau,\pi)}{\mu_w},
\]
where $\mu_g>0$ is a fixed, specified pretreatment earnings normalizer. Identify or bound both $\Pi_\tau(\pi)$ and $\Pi_\tau(\pi_1)-\Pi_\tau(\pi_0)$ for policies affecting childcare, leave, schedule flexibility, or employer incentives. In a household-firm model with declared preference, bargaining, and wage-setting assumptions, determine which separately manipulable policy mechanisms explain these differences and predict both partners' time use and employment.
\subsection*{Interpretation and scope}
An event-study divergence around childbirth is not automatically the causal effect of having a child. Fertility timing, anticipation, age profiles, and untreated counterfactual earnings require attention. Policy effects on earnings are different from effects on the birth contrast: a childcare reform can change fertility and the composition of parents. Keeping a baseline cohort and using policy-assignment effects avoids conditioning on postpolicy parenthood without justification. Household specialization can reflect preferences, constraints, taxes, or comparative advantage; observed specialization alone does not identify a unique mechanism or inefficiency. The gender earnings gap includes childless workers and prebirth differences, so it is not identical to $\Pi_\tau$. Claims about optimal policy require welfare weights, resource costs, and household preferences in addition to this earnings target.
\subsection*{Source audit and status}
[S1] analyzes professional careers, [S2] work organization, and [S3] Danish administrative event studies. They document substantial patterns and mechanisms. The displayed policy-and-birth estimands are editorial definitions; they are not claimed to be the papers' exact estimators or a named open conjecture. Cross-institutional causal mechanism identification remains a conditional research program, not a certified universal 2026 status.

\subsection*{References}
\begin{enumerate}[label=\textnormal{[S\arabic*]},leftmargin=*]
\item Marianne Bertrand, Claudia Goldin, and Lawrence F. Katz (2010). \emph{Dynamics of the Gender Gap for Young Professionals in the Financial and Corporate Sectors}. American Economic Journal: Applied Economics 2(3), 228--255. \url{https://doi.org/10.1257/app.2.3.228}. \textit{Locator:} Primary publisher, working-paper, or author record: bibliographic information and abstract; full-text theorem verification is not claimed. \textit{Role:} Career dynamics among business-school graduates. NBER working-paper title reverses Financial and Corporate; journal-title ordering supplied, without treating the two as different studies.
\item Claudia Goldin (2014). \emph{A Grand Gender Convergence: Its Last Chapter}. American Economic Review 104(4), 1091--1119. \url{https://doi.org/10.1257/aer.104.4.1091}. \textit{Locator:} Primary publisher, working-paper, or author record: bibliographic information and abstract; full-text theorem verification is not claimed. \textit{Role:} Framework connecting gender gaps and the organization of work.
\item Henrik Kleven, Camille Landais, and Jakob Egholt Søgaard (2019). \emph{Children and Gender Inequality: Evidence from Denmark}. American Economic Journal: Applied Economics 11(4). \url{https://doi.org/10.1257/app.20180010}. \textit{Locator:} Primary publisher, working-paper, or author record: bibliographic information and abstract; full-text theorem verification is not claimed. \textit{Role:} Administrative-data event-study evidence on earnings divergence around parenthood.
\end{enumerate}

\subsection*{Common conventions: Source mathematical conventions}

Unless an entry states otherwise, agent and firm sets are finite. State and action spaces are measurable, strategies and outcome maps are measurable, probability laws are specified, and expectations exist for the targets under discussion. Infinite-horizon discount factors lie in $(0,1)$ and discounted objectives are finite. Norms, tolerances and complexity input sizes must be specified for computational comparisons. Constants or primitives that appear locally are inputs, not empirical facts asserted by this catalogue.

\textbf{Identification.} A statistical model $m\in\mathcal M$ implies an observable law $P_m$ and target $\theta(m)$. At law $P$, its identified set is
\[
\Theta(P;\mathcal M)=\{\theta(m):m\in\mathcal M,\ P_m=P\}.
\]
Point identification holds when this set is a singleton, or equivalently when $P_{m_1}=P_{m_2}$ implies $\theta(m_1)=\theta(m_2)$ for all compatible models. Sharp bounds mean bounds attained, or approached when the boundary is not attained, by compatible models. Writing this set is a definition, not a solution: the substantive task is to establish informative restrictions and derive or estimate the set. An unrestricted-confounding model may give only trivial bounds.

\textbf{Inference.} Identification, estimability and valid uncertainty quantification are separate questions. Fix $\alpha\in(0,1)$. A confidence procedure $C_n$ over a declared law class $\mathcal P$ has asymptotic uniform coverage at level $1-\alpha$ when
\[
\liminf_{n\to\infty}\inf_{P\in\mathcal P}\Pr_P\{\theta(P)\in C_n\}\geq1-\alpha.
\]
For set-identified targets the coverage event must be defined separately, for example coverage of the full identified set. If an entry uses identification-indexed models instead of uniquely indexed laws, the same coverage convention is applied over those models. Pointwise limits do not establish uniform validity, and asymptotic validity does not guarantee finite-sample precision. Sample sizes, dependence structures and weak-identification sequences are part of each application.

\textbf{Counterfactuals.} $Y_i(a)$ is a potential outcome under a specified intervention, including its timing, implementation and versions. It is not $Y_i$ conditional on voluntarily choosing $a$. A full intervention vector permits interference; shorthand scalar treatment notation does not silently impose no interference. Assignment, selection, attrition, anticipation and measurement must be specified. A claim about an instrument's local effect is not automatically a population policy effect.

\textbf{Economic equilibrium.} A counterfactual model includes best responses, constraints, feasibility, market clearing, state transitions and expectations consistent with the stated information. When several equilibria exist, either a declared selector is an input or the target is set-valued. Comparisons across policy regimes must use the stated selection convention. Reduced-form relationships are not presumed invariant after an equilibrium-changing intervention.

\textbf{Optimization and welfare.} Optimization is over the stated feasible set. If attainment is not established, $\sup$ or $\inf$ and $\varepsilon$-optimal choices replace an asserted maximizer or minimizer. For example, with value $V=\sup_{a\in\mathcal A}W(a)<\infty$, an $\varepsilon$-optimal set is $\{a\in\mathcal A:W(a)\geq V-\varepsilon\}$ for $\varepsilon>0$. Benefits, costs, risk and health or environmental outcomes can be aggregated only using declared units and conversion weights. Interpersonal utility weights, the treatment of fiscal transfers and foreign profits, discounting, and the behavioral welfare criterion are explicit inputs. A comparative-static derivative additionally requires existence and appropriate differentiability of the selected counterfactual.

\textbf{Sources.} Local labels [S1], [S2], and so forth restart in every question. Locators identify the relevant abstract, model, section or result at the level actually checked; a bibliographic or abstract-level verification is not represented as inspection of an entire proof. Working papers are distinguished from later publications and changing versions. Source summaries are paraphrased. All URL checks and editorial formulations refer to the audit date stated above.
% END ECONOMICS STATEMENT
\end{document}
